\documentclass[11pt,a4paper]{article}
\usepackage[T1]{fontenc}
\usepackage[utf8]{inputenc}
\usepackage{lmodern}
\usepackage{amsmath,amssymb,amsthm,mathtools}
\usepackage{esint,graphicx}
\usepackage[a4paper,margin=22mm,headheight=14pt]{geometry}
\usepackage{microtype}
\usepackage{booktabs,longtable,array,enumitem}
\newcolumntype{L}[1]{>{\raggedright\arraybackslash}p{#1}}
\usepackage[dvipsnames]{xcolor}
\usepackage{fancyhdr}
\usepackage{hyperref}

\newtheorem{theorem}{Theorem}[section]
\newtheorem{lemma}[theorem]{Lemma}
\newtheorem{proposition}[theorem]{Proposition}
\newtheorem{corollary}[theorem]{Corollary}
\theoremstyle{remark}
\newtheorem{remark}[theorem]{Remark}
\newtheorem{definition}[theorem]{Definition}
\newtheorem{claim}[theorem]{Claim}
\setlist{nosep,leftmargin=1.5em}
\setlength{\emergencystretch}{2em}
\pagestyle{fancy}
\fancyhf{}
\fancyhead[L]{\small AIM 669}
\fancyhead[R]{\small Revised 7 November 2026}
\fancyfoot[C]{\thepage}
\renewcommand{\headrulewidth}{0.3pt}
\hypersetup{colorlinks=true,linkcolor=MidnightBlue,urlcolor=MidnightBlue,citecolor=MidnightBlue,pdftitle={AIM 669},pdfauthor={Matthew J. Colbrook; results found by AI}}
\begin{document}
\title{Metric completion need not preserve minimal Sobolev gradients}
\author{Matthew J. Colbrook\thanks{The solutions and partial results in this document were found by AI. The supplied research and the separate repository review are AI work; human peer review and formal verification are not claimed.}\\[4pt]
\small Department of Applied Mathematics and Theoretical Physics\\
\small University of Cambridge, Cambridge CB3 0WA, United Kingdom\\
\small\href{mailto:m.colbrook@damtp.cam.ac.uk}{m.colbrook@damtp.cam.ac.uk}}
\date{7 November 2026}
\maketitle
\begin{abstract}
A doubling metric measure space supporting a Poincar\'e inequality has a function with source minimal gradient at most $1$ and completion gradient $\sqrt2$. The completion complement is nonmeasurable. The conclusion concerns exactly the catalogue target; no historical priority claim is made.
\end{abstract}
\noindent\textbf{Catalogue scope.} AIM 669; complete resolution. \href{https://github.com/MColbrook/AIM/blob/1e564633494c109edfd3ec46da1366adcff925a5/problems/669-does-metric-completion-preserve-minimal-sobolev-gradients.md}{Canonical target at revision \texttt{1e56463}}.
The corrected completion-measure convention is also discussed in \cite[Section 4]{BjornBjornLahti}.

\section{Statement and argument}
\label{sec:669}

The statement in the catalog allows the original metric space to be
nonmeasurable in its completion.  Under exactly this allowance the answer is
negative.  The construction below uses the axiom of choice, but does not use
the continuum hypothesis.  The measure hypotheses are those of
\cite{BjornBjorn}: a positive complete Borel measure, finite and nonzero on
balls.  Compact inner regularity is not an additional hypothesis there.

\begin{theorem}
For every $1<p<\infty$ there is a globally doubling metric measure space $X$
supporting a global $1$-Poincar\'e inequality, whose metric and measure
completion is the Euclidean plane with Lebesgue measure, and a bounded
Lipschitz function $u$ on a ball $\Omega\subset X$ such that
\[
 g_{u,\Omega}\leq 1\quad\text{$\mu$-almost everywhere},
 \qquad
 g_{\widehat u,\Omega^\wedge}=\sqrt{2}
       \quad\text{$\mu$-almost everywhere on $\Omega$}
\]
for every Dirichlet-space extension $\widehat u$ agreeing with $u$
quasieverywhere.  Thus the constant $A_0=1$ in Problem 669 is impossible.
\end{theorem}

\subsection{A set excluding oblique curves}

\begin{lemma}
There is a set $A\subset\mathbb R$, containing $0$, such that both $A$ and
$\mathbb R\setminus A$ meet every positive-measure Borel subset of
$\mathbb R$, and every rectifiable curve in
\[
 X=(A\times\mathbb R)\cup(\mathbb R\times A)
\]
has an axis-parallel tangent at almost every arclength parameter.
\end{lemma}

\begin{proof}
Call a Lipschitz map $\gamma:[0,1]\to\mathbb R^2$ oblique if
\[
 \mathcal L^1\{t:\gamma_1'(t)\ne0\text{ and }\gamma_2'(t)\ne0\}>0.
\]
There are at most $\mathfrak c$ such maps, since continuous maps are
determined by their values on a countable dense set.  There are exactly
$\mathfrak c$ positive-measure Borel subsets of $\mathbb R$.

For each oblique map there is a positive-measure set $S\subset[0,1]$ on
which both coordinate functions are injective.  Indeed, choose $\delta>0$
so that the set of differentiability points satisfying
$|\gamma_i'(t)|\ge\delta$, $i=1,2$, has positive measure.  At each such
point, for all sufficiently small nonzero $h$ with $t+h\in[0,1]$,
\[
 |\gamma_i(t+h)-\gamma_i(t)|\ge\tfrac{\delta}{2}|h|,
 \qquad i=1,2.
\]
The measurable subsets on which this holds for every $0<|h|<1/m$,
$m\in\mathbb N$, cover that set.  One of them has positive measure.
Intersect it with an interval of length less than $1/m$ having
positive-measure intersection.  The resulting set $S$ satisfies the
displayed inequality for every pair of its points.  In particular each
coordinate is injective on $S$, and $|S|=\mathfrak c$.

Well-order, in one list of cardinality $\mathfrak c$, all the oblique maps
and all positive-measure Borel subsets of $\mathbb R$.  Maintain disjoint
sets of real numbers assigned to $A$ and to $B$, initially assigning $0$
to $A$.  At a Borel-set task choose two previously unassigned points of
that set, assigning one to $A$ and one to $B$.  At an oblique-map task
choose $t\in S$ such that both $\gamma_1(t)$ and $\gamma_2(t)$ are
previously unassigned, and assign both to $B$.

These choices are always possible.  At a stage indexed by an ordinal
$\alpha<\mathfrak c$, fewer than $\mathfrak c$ numbers have been assigned.
Every positive-measure Borel set has cardinality $\mathfrak c$.  At a
curve task, coordinate injectivity on $S$ implies that fewer than
$\mathfrak c$ parameters are prohibited by earlier assignments.  This
argument does not require $\mathfrak c$ to be a regular cardinal.

After all tasks assign every remaining number to $B$, and let $A$ be the
set assigned to $A$.  The Borel-set tasks give the first assertion.  Every
oblique map has a selected point in $B\times B$, so cannot have its whole
image in $X$.  Apply this observation to arclength parametrizations of
rectifiable curves, affinely rescaled to $[0,1]$, to obtain the second
assertion.
\end{proof}

\subsection{The trace measure}

Equip $X$ with the Euclidean subspace metric.  The set $A$ is dense, so
$\widehat X=\mathbb R^2$.  Every positive-measure ambient Borel set meets
$A\times\mathbb R$, hence meets $X$: by Fubini, the set of first
coordinates with positive-measure vertical section has positive measure
and therefore meets $A$.  Thus $\mathbb R^2\setminus X$ has inner
Lebesgue measure zero.

Define a measure on the trace Borel sigma algebra by
\begin{equation}
 \mu(E\cap X)=\mathcal L^2(E),\qquad E\subset\mathbb R^2\text{ Borel},
 \label{eq:669trace}
\end{equation}
and complete it.  This is well-defined: if two ambient Borel sets have the
same trace, their Borel symmetric difference is contained in
$\mathbb R^2\setminus X$ and is null.  If traces are pairwise disjoint,
their representatives have null pairwise intersections; disjointifying
the representatives proves countable additivity.  Every ball satisfies
\[
 \mu(B_X(c,r))=\pi r^2,
 \qquad \mu(B_X(c,2r))=4\mu(B_X(c,r)).
\]
The completion measure prescribed in the question is consequently
ordinary completed Lebesgue measure.

We will use two elementary properties of this trace construction.  Every
$\mu$-null set is contained in $N\cap X$ for some ambient null Borel set
$N$.  Also, every $\mu$-measurable real-valued function has an ambient
Lebesgue-measurable representative agreeing with it outside such a trace
null set.  To see the latter, first take a Borel version for the completed
trace measure.  A real-valued Borel function on a subspace has an ambient
Borel extension: construct it from Borel extensions of simple-function
approximations and their pointwise limits.  Trace integrals then equal
ambient integrals by \eqref{eq:669trace} and simple-function approximation.
These are also the measure-extension conventions in the corrigendum to
\cite{BjornBjorn}.

\subsection{A global Poincar\'e inequality}

We prove the $1$-Poincar\'e inequality directly; H\"older's inequality
then gives the $p$-Poincar\'e inequality for every $p\ge1$.
Let $B=B_X(c,r)$, let $v$ be integrable on $2B$, and let $g$ be a Borel
upper gradient of $v$ there.  We may suppose $g\in L^1(2B)$, since
otherwise the required inequality is trivial.  Choose ambient Borel
representatives $V,G$ on the Euclidean ball $\widehat B(c,2r)$, with
$G\ge0$, and one ambient null Borel set $N$ outside whose trace both
representatives agree with the original functions.  Let
$Q=I_1\times I_2$ be the open axis-aligned square of side $2r$ centered
at $c$.  Then
\[
 \widehat B(c,r)\subset Q\subset\widehat B(c,2r).
\]

By Fubini, outside a null set $Z\subset I_2$ the functions
$V(\cdot,y)$ and $G(\cdot,y)$ are integrable on $I_1$ and the section
$N_y$ is null.  For every $y\in(A\cap I_2)\setminus Z$, the whole
horizontal chord $I_1\times\{y\}$ lies in $X$.  Thus, for every pair
$s,t\in I_1\setminus N_y$, the upper-gradient inequality gives
\begin{equation}
 |V(s,y)-V(t,y)|
       \le\int_{\min(s,t)}^{\max(s,t)}G(x,y)\,dx.
 \label{eq:669line}
\end{equation}
This gives an absolutely continuous representative of the slice and
$|\partial_1V(\cdot,y)|\le G(\cdot,y)$ almost everywhere.  For example,
absolute continuity of the integral on the right first extends $V$ from
the full-measure set $I_1\setminus N_y$ continuously to $I_1$; applying
\eqref{eq:669line} to finitely many disjoint intervals proves absolute
continuity and the derivative bound.

The passage from the possibly nonmeasurable family $A$ to almost every
line is justified as follows.  Fix a countable family
$\mathcal T\subset C_c^\infty(I_1)$ dense in the $C^1$ topology on each
of a countable compact exhaustion of $I_1$.  The set of parameters
$y\in I_2\setminus Z$ for which at least one inequality
\begin{equation}
 \left|\int_{I_1}V(x,y)\varphi'(x)\,dx\right|
       \le\int_{I_1}G(x,y)|\varphi(x)|\,dx,
       \qquad \varphi\in\mathcal T,
 \label{eq:669tests}
\end{equation}
fails is Lebesgue measurable, by Fubini and countability.  By the previous
paragraph it is disjoint from $A$.  It must therefore be null: every
positive-measure measurable set contains a positive-measure Borel set,
and $A$ meets all such sets.  Density extends
\eqref{eq:669tests} to every smooth compactly supported test function.
The one-dimensional distributional characterization of absolute
continuity gives $|\partial_1V|\le G$ on almost every horizontal slice.
Repeating the argument on the vertical lines whose first coordinate is
in $A$ gives $|\partial_2V|\le G$ on almost every vertical slice.

Fubini and integration by parts on the slices now show that
$V\in W^{1,1}(Q)$ and
\[
 |\nabla V|\le\sqrt2\,G\quad\text{almost everywhere on }Q.
\]
The usual Euclidean $L^1$ Poincar\'e inequality on $\widehat B(c,r)$,
followed by the ball-volume ratio, therefore gives
\[
 \fint_B|v-v_B|\,d\mu
 \le C r\fint_{\widehat B(c,r)}|\nabla V|\,d\mathcal L^2
 \le C' r\fint_{2B}g\,d\mu.
\]
The constants are absolute and independent of $B,v,g$.  This proves the
claimed global Poincar\'e inequality with dilation factor $2$.

\subsection{The gradient gap}

Take $B_0=B_X(0,2)$ and $\Omega=B_X(0,1)$.  The definition
$\Omega^\wedge=\widehat X\setminus\overline{X\setminus\Omega}$ gives
the ordinary open Euclidean unit disk, since $A$ is dense.  Set
$u(x_1,x_2)=x_1+x_2$ on $\Omega$.  Every rectifiable curve in $\Omega$,
parametrized by arclength, has at most one nonzero coordinate derivative
almost everywhere.  Consequently
\[
 |(u\circ\gamma)'|
       =|\gamma_1'+\gamma_2'|=|\gamma'|
       \quad\text{almost everywhere}.
\]
Thus $g=1$ is an upper gradient of $u$, belongs to $L^p(\Omega)$, and
$g_{u,\Omega}\le1$ almost everywhere.

Suppose $\widehat u\in D^p(\Omega^\wedge)$ agrees with $u$ outside a
set of zero $C_p^X$ capacity.  Zero Sobolev capacity implies zero
$\mu$-measure, since the capacity includes the $L^p$ term.  The
disagreement set is therefore contained in $N\cap X$ for an ambient
null Borel set $N$.  Since $X$ meets every positive-measure ambient Borel
set, the measurable function $\widehat u-(x_1+x_2)$ must vanish
Lebesgue-almost everywhere on $\Omega^\wedge$.  Indeed, a positive-measure
set of disagreement outside $N$ would contain a positive-measure Borel
set meeting $X$, a contradiction.

Euclidean Sobolev theory identifies the minimal $p$-weak upper gradient
with the length of the weak gradient.  It follows that
\[
 g_{\widehat u,\Omega^\wedge}
       =|\nabla(x_1+x_2)|=\sqrt2
       \quad\text{Lebesgue-almost everywhere}.
\]
By the trace-measure identity this also holds $\mu$-almost everywhere on
$\Omega$.  This proves the theorem.  In this example any valid comparison
constant $A_0$ is at least $\sqrt2$.

The construction addresses the exact generality of the catalog.  It
does not assert failure when $X$ is required to be measurable in its
completion.

\begingroup\small
\begin{thebibliography}{99}
\bibitem{BjornBjornLahti} A. Bj\"orn, J. Bj\"orn and P. Lahti, \emph{Removable sets for Newtonian Sobolev spaces and a characterization of $p$-path almost open sets}, Rev. Mat. Iberoam. \textbf{39} (2023), 1143--1180. \url{https://arxiv.org/abs/2105.09012v2}. DOI: \url{https://doi.org/10.4171/RMI/1419}.
\bibitem{BjornBjorn} A. Bj\"orn and J. Bj\"orn, \emph{Poincar\'e inequalities and Newtonian Sobolev functions on noncomplete metric spaces}, J. Differential Equations \textbf{266} (2019), 44--69; corrigendum, \textbf{285} (2021), 493--495. \url{https://arxiv.org/abs/1705.02253}; \url{https://doi.org/10.1016/j.jde.2021.03.008}. DOI: \url{https://doi.org/10.1016/j.jde.2018.07.029}.
\end{thebibliography}
\endgroup
\end{document}
